% !TEX root = ../main.tex
\chapter{Glossário}
\label{ap:glossario}

\begin{description}[style=nextline,font=\bfseries]
  \item[SAPHO] \emph{Scalable-Architecture Processor for Hardware Optimization}: o processador \emph{soft-core} e, por extensão, a plataforma completa (AURORA + YANC + toolchain) e o canal de distribuição.
  \item[AURORA] \emph{Advanced Utility Running Optimized Resource Architectures}: a IDE de desktop da plataforma.
  \item[YANC] \emph{Yet Another Compiler}: a suíte de compiladores que traduz \cmm{}/C++ em processador Verilog.
  \item[\cmm{} (CMM)] A linguagem de programação da plataforma; dialeto de C com complexos nativos e notação de Dirac. Arquivos \file{.cmm}.
  \item[PRISM] \emph{Processor Rendering Interface for Schematic Models}: o visualizador de RTL da AURORA.
  \item[Aurora Intelligence] A assistente de IA integrada à AURORA.
  \item[Dagr] O painel de controle de versão (Git) da AURORA.
  \item[NIPS-CERN] O Núcleo de Instrumentação e Processamento de Sinais da UFJF, em colaboração com o CERN --- o laboratório que desenvolve a plataforma.
  \item[soft-processor / soft-core] Processador descrito em HDL e implementado na malha reconfigurável de um FPGA.
  \item[FPGA] \emph{Field-Programmable Gate Array}: circuito integrado reconfigurável.
  \item[Verilog / HDL] Linguagem de descrição de hardware usada pela plataforma.
  \item[\file{.spf}] \emph{SAPHO Project File}: o arquivo JSON que define um projeto.
  \item[\file{.asm}] O assembly intermediário do SAPHO, entre o compilador e o montador.
  \item[\file{.mif}] \emph{Memory Initialization File}: imagem de memória (programa ou dados) embutida no design.
  \item[testbench] O banco de testes da simulação: um módulo Verilog (\file{_tb.v}) ou um script Python (cocotb) que estimula o processador.
  \item[top-level] O módulo Verilog que integra os processadores e blocos do projeto; topo da síntese.
  \item[ULA] Unidade Lógico-Aritmética: o conjunto de operadores do processador; no SAPHO, só os usados são instanciados.
  \item[acumulador (ACC)] O registrador central do processador SAPHO, destino de toda operação da ULA.
  \item[Harvard] Arquitetura com memórias de programa e dados separadas.
  \item[opcode / ISA] Código de operação de uma instrução / o conjunto de instruções do processador.
  \item[Icarus Verilog / vvp] O simulador de Verilog padrão (interpreta o design).
  \item[Verilator] Simulador de alto desempenho que converte o design em executável C++.
  \item[cocotb] \emph{Framework} Python para escrever \emph{testbenches}.
  \item[Yosys] Ferramenta de síntese usada pelo PRISM e pela visão Hierarquia.
  \item[GTKWave / Surfer] Os dois visualizadores de formas de onda da plataforma.
  \item[VCD / FST] Formatos de arquivo de ondas gerados pela simulação (textual / binário compacto).
  \item[\file{.gtkw} / \file{.surf.ron}] Arquivos de layout de ondas (GTKWave / Surfer): quais sinais, cores e grupos mostrar.
  \item[Monaco] O motor de edição de código (o mesmo do VS Code).
  \item[LSP] \emph{Language Server Protocol}: o mecanismo dos diagnósticos e navegação de código (Verible, slang).
  \item[BYOK] \emph{Bring Your Own Key}: o modelo de chaves da Aurora Intelligence.
  \item[MCP] \emph{Model Context Protocol}: como as CLIs de IA acessam as ferramentas da AURORA.
  \item[bit-reverso] Endereçamento com bits do índice invertidos (\code{x[k)}), usado em FFT.
  \item[DLP] Dispositivos Lógicos Programáveis: a disciplina da UFJF que usa a plataforma no ensino.
\end{description}
